% !TEX program = pdflatex
% IEEE 会议论文模板（DAC/ICCAD 等 EDA 方向投稿）
% IEEEtran.cls 已随 MiKTeX 安装
% 编译: latexmk -pdf main.tex（或 bin/latex_build.py build --dir . --engine pdflatex）
\documentclass[conference]{IEEEtran}
\usepackage{cite}
\usepackage{amsmath,amssymb,amsfonts}
\usepackage{algorithmic}
\usepackage{graphicx}
\usepackage{textcomp}
\usepackage{xcolor}
\def\BibTeX{{\rm B\kern-.05em{\sc i\kern-.025em b}\kern-.08em
    T\kern-.1667em\lower.7ex\hbox{E}\kern-.125emX}}

\begin{document}

\title{Paper Title: A Shortened Version}

\author{\IEEEauthorblockN{Author One}
\IEEEauthorblockA{\textit{School of Electronic Information and Electrical Engineering}\\
Shanghai Jiao Tong University, Shanghai, China \\
author@example.com}
\and
\IEEEauthorblockN{Author Two}
\IEEEauthorblockA{\textit{Department of XXX}\\
Shanghai Jiao Tong University, Shanghai, China \\
author2@example.com}}

\maketitle

\begin{abstract}
This is the abstract. Describe the problem, method, and result.
\end{abstract}

\begin{IEEEkeywords}
circuit, LLM, EDA
\end{IEEEkeywords}

\section{Introduction}
Background \cite{dac2026} and contribution bullets:
\begin{itemize}
\item Contribution one.
\item Contribution two.
\end{itemize}

\section{Method}
See Eq.~\eqref{eq:loss}.
\begin{equation}
\label{eq:loss}
  \mathcal{L} = \mathbb{E}[-\log p_\theta(y \mid x)]
\end{equation}

\section{Experiments}

\section{Conclusion}

\section*{Acknowledgment}

\bibliographystyle{IEEEtran}
\bibliography{refs}

\end{document}
